\DocumentMetadata{
  pdfversion=2.0,
  pdfstandard=ua-2,
  lang=en-US,
  tagging=on,
  tagging-setup={math/setup=mathml-SE,tabsorder=structure}
}
\documentclass[11pt,letterpaper]{article}
\usepackage[margin=0.68in,includefoot]{geometry}
\usepackage{newcomputermodern}
\usepackage{amsmath,amscd,mathtools}
\usepackage{tikz,adjustbox}
\providecommand{\square}{\mdlgwhtsquare}
\usepackage[hidelinks]{hyperref}
\hypersetup{
  pdftitle={Topology First-Year Exam, August 2019, Part 1},
  pdfauthor={University of Florida Department of Mathematics},
  pdfsubject={Graduate mathematics examination},
  pdfdisplaydoctitle=true,
  bookmarksopen=true
}
\setcounter{secnumdepth}{0}
\setlength{\parindent}{0pt}
\setlength{\parskip}{0.28em}
\setlength{\emergencystretch}{3em}
\setlength{\leftmargini}{2.15em}
\setlength{\leftmarginii}{2.65em}
\setlength{\leftmarginiii}{3em}
\pagestyle{empty}
\raggedbottom
\allowdisplaybreaks
\NewTaggingSocketPlug{block/list/label}{examproblem}{%
  \tagstructbegin{tag=Lbl}\tagstructend%
  \tagstructbegin{tag=\LBody}%
  \tagstructbegin{tag=\ProblemHeadingTag,title-o={\ProblemHeadingText}}%
  \tagmcbegin{tag=\ProblemHeadingTag,actualtext={\ProblemHeadingText}}%
  #2%
  \tagmcend%
  \tagstructend%
}
\newenvironment{examproblems}[2]{%
  \def\ProblemHeadingTag{#2}%
  \AssignTaggingSocketPlug{block/list/label}{examproblem}%
  \begin{enumerate}%
  \setcounter{enumi}{#1}%
  \renewcommand{\labelenumi}{\textbf{\theenumi.}}%
  \setlength{\topsep}{0.35em}%
  \setlength{\partopsep}{0pt}%
  \setlength{\itemsep}{0.48em}%
  \setlength{\parsep}{0.18em}%
}{\end{enumerate}}
\newenvironment{examsubparts}{%
  \AssignTaggingSocketPlug{block/list/label}{default}%
  \begin{enumerate}%
  \setlength{\topsep}{0.18em}%
  \setlength{\partopsep}{0pt}%
  \setlength{\itemsep}{0.16em}%
  \setlength{\parsep}{0.1em}%
}{\end{enumerate}}
\newenvironment{examinstructions}{%
  \par\begingroup\itshape
}{%
  \par\endgroup\vspace{0.18em}
}
\makeatletter
\renewcommand\subsection{\@startsection{subsection}{2}{0pt}%
  {-1.25ex plus -.25ex minus -.1ex}%
  {0.55ex plus .1ex}%
  {\normalfont\large\bfseries}}
\renewcommand{\@maketitle}{%
  \newpage\null\vskip -1em%
  {\footnotesize\bfseries
    \href{https://www.ufl.edu/}{University of Florida}, \href{https://math.ufl.edu/}{Department of Mathematics}\par}%
  \vskip -0.5em%
  \tagstructbegin{tag=H1,title={Topology First-Year Exam, August 2019, Part 1}}%
  \tagpdfparaOff%
  \tagmcbegin{tag=H1}%
    {\LARGE\bfseries\href{https://gma.math.ufl.edu/exams/topology-first-year/}{Topology First-Year Exam}, August 2019, Part 1\par}%
  \tagmcend%
  \tagpdfparaOn%
  \tagstructend%
  \vskip 0.25em%
  \tagmcbegin{artifact}%
    \hrule height 1pt%
  \tagmcend%
  \vskip 1em%
}
\makeatother
\title{Topology First-Year Exam, August 2019, Part 1}
\author{}
\date{}
\begin{document}
\maketitle
\thispagestyle{empty}
\begin{examinstructions}
Work all problems.
\end{examinstructions}
\subsection{Theorems and Proofs.}
\begin{examinstructions}
In this section state and prove each theorem. The proofs should be complete without being tedious.
\end{examinstructions}
\begin{examproblems}{0}{H3}
\def\ProblemHeadingText{Problem 1.}
\item
Let~\(X\) be a metric space and suppose that \(A_\lambda \subset X\) is connected for all \(\lambda \in \Lambda\). Suppose that there is an \(x_0 \in X\) such that \(x_0 \in A_\lambda\) for all~\(\lambda\). Show that \(\bigcup_{\lambda \in \Lambda} A_\lambda\) is connected.
\def\ProblemHeadingText{Problem 2.}
\item
State and prove the Contraction Mapping Theorem.
\def\ProblemHeadingText{Problem 3.}
\item
Let~\(X\) be a Hausdorff space. Show that if \(A \subset X\) is compact, then~\(A\) is closed.
\def\ProblemHeadingText{Problem 4.}
\item
Let~\(X\) be a topological space. Suppose that \(f:X \to Y\) is continuous. Show that if \(A \subset X\) is connected, then \(f(A) \subset Y\) is connected.
\def\ProblemHeadingText{Problem 5.}
\item
Suppose that~\(X\) and~\(Y\) are metric spaces and that~\(X\) is compact. Show that if \(f:X \to Y\) is continuous, then~\(f\) is uniformly continuous.
\end{examproblems}
\subsection{Examples and Minor Proofs.}
\begin{examinstructions}
In this section, give brief examples, counterexamples, or quick proofs.
\end{examinstructions}
\begin{examproblems}{5}{H3}
\def\ProblemHeadingText{Problem 6.}
\item
State Urysohn’s Lemma.
\def\ProblemHeadingText{Problem 7.}
\item
Is it possible for there to be a continuous function \(f:[0,1] \to \{0,1\}\) which is onto?
\def\ProblemHeadingText{Problem 8.}
\item
State the Bolzano–Weierstrass Theorem.
\def\ProblemHeadingText{Problem 9.}
\item
Define a continuous function \(f:C \to [0,1]\) which is onto where~\(C\) is the Cantor set and \([0,1]\) is the unit interval.
\def\ProblemHeadingText{Problem 10.}
\item
Give an example of a space~\(X\) which is \(T_1\) but not Hausdorff.
\def\ProblemHeadingText{Problem 11.}
\item
State the Tietze Extension Theorem.
\def\ProblemHeadingText{Problem 12.}
\item
Define a quotient map.
\def\ProblemHeadingText{Problem 13.}
\item
State the Baire Category Theorem.
\def\ProblemHeadingText{Problem 14.}
\item
What does it mean for a topological space~\(X\) to be Hausdorff? Regular? Normal?
\def\ProblemHeadingText{Problem 15.}
\item
What does it mean for a space~\(X\) to be first countable? Show that every metric space is first-countable.
\end{examproblems}
\end{document}
